% Standalone problem document, generated from the adjacent README.md.
% Compile with XeLaTeX. Mathematical content is maintained in Markdown.
\documentclass[11pt,a4paper]{article}
\usepackage[fontsize=10.5pt]{scrextend}
\usepackage[margin=22mm,headheight=15pt,headsep=9mm,footskip=11mm]{geometry}
\usepackage{fontspec}
\setmainfont{texgyrepagella}[Extension=.otf,UprightFont=*-regular,BoldFont=*-bold,ItalicFont=*-italic,BoldItalicFont=*-bolditalic]
\setsansfont{texgyreheros}[Extension=.otf,UprightFont=*-regular,BoldFont=*-bold,ItalicFont=*-italic,BoldItalicFont=*-bolditalic]
\setmonofont{texgyrecursor}[Extension=.otf,UprightFont=*-regular,BoldFont=*-bold,ItalicFont=*-italic,BoldItalicFont=*-bolditalic,Scale=MatchLowercase]
\usepackage{amsmath,amssymb}
\usepackage{unicode-math}
\setmathfont{texgyrepagella-math.otf}
\usepackage{xcolor,microtype,enumitem,needspace,fancyhdr}
\usepackage{bookmark}
\definecolor{ink}{HTML}{17344B}
\definecolor{accent}{HTML}{197D80}
\definecolor{muted}{HTML}{596773}
\hypersetup{colorlinks=true,linkcolor=accent,urlcolor=accent,pdftitle={AC-01 - Is
the matrix multiplication exponent
two?},pdfauthor={Open Problems in Numerical Linear Algebra}}
\urlstyle{same}
\setlength{\parindent}{0pt}
\setlength{\parskip}{5pt plus 1pt minus 1pt}
\linespread{1.04}
\setlength{\emergencystretch}{3em}
\widowpenalty=10000
\clubpenalty=10000
\displaywidowpenalty=10000
\allowdisplaybreaks[1]
\setlist{nosep,leftmargin=1.5em}
\providecommand{\tightlist}{\setlength{\itemsep}{0pt}\setlength{\parskip}{0pt}}
\providecommand{\pandocbounded}[1]{#1}
\pagestyle{fancy}
\fancyhf{}
\fancyhead[L]{\sffamily\footnotesize\color{muted}OPEN PROBLEMS IN NUMERICAL LINEAR ALGEBRA}
\fancyhead[R]{\sffamily\bfseries\footnotesize\color{accent}AC-01}
\fancyfoot[L]{\parbox[b]{0.9\textwidth}{\sffamily\fontsize{8}{10}\selectfont\color{muted}Editorial ratings; a bounded search cannot certify that a problem remains unsolved.\\Literature check: 2026-09-10}}
\fancyfoot[R]{\sffamily\footnotesize\color{muted}\thepage}
\renewcommand{\headrulewidth}{0.3pt}
\renewcommand{\headrule}{\hbox to\headwidth{\color{accent}\leaders\hrule height \headrulewidth\hfill}}
\makeatletter
\renewcommand{\section}{\Needspace{5\baselineskip}\@startsection{section}{1}{0pt}{12pt plus 2pt minus 2pt}{4pt}{\sffamily\bfseries\large\color{ink}}}
\renewcommand{\subsection}{\Needspace{4\baselineskip}\@startsection{subsection}{2}{0pt}{9pt plus 1pt minus 1pt}{3pt}{\sffamily\bfseries\normalsize\color{ink}}}
\makeatother
\setcounter{secnumdepth}{0}
\begin{document}
{\sffamily\small\color{muted}Arithmetic and complexity\par}
\vspace{3pt}
{\raggedright\hyphenpenalty=10000\exhyphenpenalty=10000\sffamily\bfseries\fontsize{21}{24}\selectfont\color{ink}Is
the matrix multiplication exponent two?\par}
\vspace{7pt}
{\sffamily\small\textbf{Difficulty:} extreme\quad\textbf{Importance:} broadly
interesting\par}
{\sffamily\small\bfseries\color{accent}Status: Open\par}
\vspace{4pt}
{\color{accent}\hrule height 0.8pt}
\vspace{6pt}
\textbf{Rating rationale:} Extreme because closing the exponent gap is a
longstanding central barrier in algebraic algorithms; broad importance
follows from matrix multiplication's role throughout NLA and
computational complexity.\\
\textbf{Topic:} arithmetic complexity of dense matrix multiplication

\subsection{Problem statement}\label{problem-statement}

Over the field \(\mathbb C\), let \(M(n)\) be the minimum number of
scalar additions, subtractions, multiplications, and divisions in an
arithmetic straight-line program computing every entry of \(AB\) for
arbitrary \(A,B\in\mathbb C^{n\times n}\). Programs must be defined on
their intended inputs. Set \(\omega=\inf\{\tau:M(n)=O(n^\tau)\}\). Is
\(\omega=2\)? Equivalently, for every \(\varepsilon>0\), can these
products be computed in \(O(n^{2+\varepsilon})\) arithmetic operations?
This asks about asymptotic arithmetic cost; it does not assert an
\(O(n^2)\) algorithm or a floating-point stability guarantee.

\subsection{Why it matters}\label{why-it-matters}

Matrix multiplication is a basic cost driver for dense NLA.

\subsection{References and status}\label{references-and-status}

M. Bläser,
\href{https://theoryofcomputing.org/articles/gs005/gs005.pdf}{\emph{Fast
Matrix Multiplication}}, Graduate Surveys 5 (2013), §§1, 5, provides the
computational model. E. Dupont et al.,
\href{https://arxiv.org/abs/2608.16884}{\emph{Improving the matrix
multiplication exponent with modern optimization and AlphaEvolve}}
(2026), abstract and §1, reports \(\omega<2.371177\), which does not
reach two. Searches for ``matrix multiplication exponent 2026'' and
``omega equals 2 proof'' located improvements, not a resolution.
\textbf{Admitted: no resolution located.}

\subsection{Status check --- 2026-09-10}\label{status-check-2026-09-10}

Rechecked \href{https://arxiv.org/abs/2608.16884}{Dupont et al., August
2026}, and searched for exponent-two proofs and newer
matrix-multiplication bounds. Its reported bound is still strictly above
two, at 2.371177. No proof of exponent two or a strict lower bound above
two was located. Faster finite-size identities and improved numerical
optimizations of existing bounds do not decide the asymptotic equality.
\end{document}
